% Pregunta 1. Perfil cíclico (bloque I) · peso alto
\section*{Pregunta 1. Perfil cíclico}

\begin{enunciado}
\textbf{1.} Describa brevemente el \textbf{perfil cíclico} de la economía española a lo largo del periodo 2007-2025, desde la óptica de la demanda y de la oferta, identificando sus distintas fases, y valore su situación coyuntural a la luz de los datos trimestrales de 2026 y de las previsiones para 2026 y 2027. ¿Qué indica la brecha de producción sobre la posición cíclica? ¿A qué podría deberse que, desde 2019, el PIB por habitante haya crecido bastante menos que el PIB?

\smallskip
\textbf{Peso alto} · \tiempo{23} · guía: I.2 (fases), I.3 (medidores), I.4 (demanda), I.5 (oferta), I.6 (diagnóstico) · cuadros 1, 2 y 2~bis; gráficos 1 y 2.
\end{enunciado}

\begin{examen}
\textit{La economía española ha recorrido entre 2007 y 2025 un ciclo completo con 2 recesiones y encadena desde 2021 una expansión más intensa que la de la zona del euro, apoyada en la demanda interna, el empleo y el aumento de la población. Produce ya por encima de su potencial y mantiene superávit exterior, pero crece más por acumular trabajo que por ganar productividad, y afronta en 2026 un nuevo choque de precios energéticos.}

\rotulo{a) Las fases del ciclo, por el lado de la demanda y de la oferta}
El cuadro recoge más de un ciclo: el de la crisis financiera y de la deuda soberana y el de la pandemia. Con el crecimiento del PIB real y la brecha de producción (cuadro~1) se distinguen 5 etapas:

\textbf{1.}~\textbf{Final de la expansión (2007-2008).} El PIB pasa de crecer un 3,5~\% a un 0,8~\% con la brecha todavía positiva (+3,3~\% en 2007): la economía producía por encima de su potencial, con una necesidad de financiación exterior del 9,0~\% del PIB (cuadro~7) que anticipaba un ajuste.


\textbf{2.}~\textbf{Doble recesión (2009-2013).} Una primera caída en 2009 ($-3{,}8$~\%), un estancamiento en 2010-2011 y una segunda recesión en 2012-2013 ($-2{,}9$ y $-1{,}4$~\%): una recesión en forma de W, con la brecha en un mínimo del $-9{,}6$~\% en 2013 y el paro en el 26,1~\%. La primera caída responde al fin del auge inmobiliario y del crédito; la segunda coincide con la crisis de la deuda soberana (diferencial de 435~pb en 2012, cuadro~9) y con la consolidación fiscal.


\textbf{3.}~\textbf{Recuperación y expansión (2014-2019).} El PIB crece un 2,6~\% de media; la brecha se cierra en 2017 y llega al +3,1~\% en 2019.


\textbf{4.}~\textbf{Pandemia (2020).} Una caída del 10,9~\%, mucho mayor que la de la zona del euro (en torno al 6~\%, gráfico~2), por el peso de los servicios de contacto y del turismo.


\textbf{5.}~\textbf{Rebote y nueva expansión (2021-2025).} Tras crecer un 6,7 y un 6,4~\% en 2021-2022, el PIB mantiene ritmos del 2,4-3,7~\% y la brecha es positiva desde 2022 (+2,0~\% en 2025).


Por el lado de la \textbf{demanda}, la nacional es \emph{procíclica} y la externa \emph{contracíclica}. La demanda nacional explica tanto las expansiones (4,2~pp en 2007 y en 2015; 3,4~pp en 2025) como las recesiones ($-6{,}5$~pp en 2009 y $-4{,}9$~pp en 2012), mientras la externa amortigua las caídas (+2,8~pp en 2009 y +2,0~pp en 2012), más por el desplome de las importaciones ($-18{,}3$~\% en 2009) que por el empuje de las exportaciones. Dentro de la nacional, la inversión es, como es habitual, el componente más volátil (la formación bruta de capital fijo cae un 17,1~\% en 2009 y la inversión en vivienda, un 20,2~\%), y el consumo privado sigue al PIB con menos amplitud, salvo en 2020 ($-12{,}1$~\%), cuando el cierre de la actividad lo convirtió en el canal principal del choque. El consumo público amortigua la primera recesión (+4,9 y +3,8~\% en 2008-2009) pero se recorta en plena segunda ($-3{,}9$~\% en 2012), un indicio de política fiscal procíclica que se retoma en la pregunta~9. La demanda externa vuelve a sumar en 2022 (+2,3~pp), con la vuelta del turismo, y resta en 2025 ($-0{,}8$~pp) porque las importaciones (+5,5~\%) crecen más que las exportaciones (+2,6~\%) al tirar la demanda interna: una expansión, por tanto, guiada desde dentro.

Por el lado de la \textbf{oferta} (cuadro~2), la expansión de 2014-2019 fue generalizada; la actual es sobre todo de servicios de mercado, con las actividades profesionales (+8,3 y +7,2~\% en 2024-2025) y la información y las comunicaciones a la cabeza, mientras la industria apenas avanza ($-2{,}1$~\% en 2023 y +0,6~\% en 2025), lo que sería coherente con el encarecimiento de la energía y la pérdida de competitividad en costes (pregunta~4). La agricultura es volátil y acíclica ($-16{,}9$~\% en 2022, probablemente por la sequía, y +15,0 y +11,9~\% después), y la construcción, que fue la gran perdedora de 2009-2013, repunta en 2025 (+6,0~\%).

\rotulo{b) La situación en 2026 y las previsiones}
Los datos trimestrales confirman que la expansión continúa: el PIB crece un 0,6 y un 0,7~\% en el primer y el segundo trimestre de 2026 (gráfico~1), un ritmo anualizado del 2,5-2,8~\%, apoyado en el consumo privado (0,67~pp en el segundo trimestre), mientras la demanda externa vuelve a restar ($-0{,}46$~pp). Con avances del 0,4~\% en la segunda mitad del año, el crecimiento medio de 2026 rondaría el 2,5~\%, en línea con la previsión de la Comisión (2,4~\%). Las previsiones (2,4~\% en 2026 y 1,9~\% en 2027) dibujan una \textbf{desaceleración suave}: el PIB crecería menos que su potencial (2,7 y 2,5~\%) y la brecha se reduciría al 1,7 y al 1,2~\%. Es un aterrizaje suave plausible, pero con riesgos a la baja (el petróleo, la subida de tipos del BCE y los aranceles) que se analizan en la pregunta~10.

\rotulo{c) Qué dice la brecha de producción}
La brecha es positiva desde 2022 y está en torno al 2~\% del potencial: \textbf{la economía produce por encima de su potencial}, en una fase de expansión madura. Es coherente con lo que se verá en otros bloques: una inflación por encima de la de la zona del euro, unos costes laborales unitarios que crecen el 4,4~\% y un mercado de trabajo con más vacantes que antes de la pandemia. Sin embargo, la brecha apenas se amplía desde 2024 (del 1,9 al 2,0~\%): el PIB crece ya al ritmo del potencial, que la Comisión estima en el 2,7~\%, muy por encima del $-0{,}1$-1,0~\% de 2014-2019. \uline{Un potencial tan alto para una economía avanzada no se explica por la productividad, sino por la población}: la economía crece con fuerza sin que la brecha se abra porque llegan trabajadores. Conviene leerla con prudencia: la brecha es una estimación, se revisa mucho y es especialmente incierta cuando la población activa cambia deprisa.

\rotulo{d) Por qué el PIB por habitante crece menos que el PIB}
Entre 2019 y 2025 el PIB crece un 10,0~\% y el PIB por habitante solo un 5,0~\% (gráfico~2), porque la población aumenta casi un 5~\% (cuadro~3), impulsada por la inmigración: la población extranjera pasa del 10,3 al 14,1~\% del total. La descomposición del cuadro~2~bis lo precisa: de los 1,59~puntos de crecimiento medio anual del PIB, 0,79 los aporta la población y 0,95 la mayor tasa de ocupación, mientras las horas por ocupado restan 0,46 y la productividad por hora solo suma 0,33. Es un \textbf{crecimiento extensivo}, apoyado en más personas trabajando y no en que cada hora produzca más. En la zona del euro la diferencia entre PIB y PIB por habitante es mucho menor (6,8 y 4,7~\%), y desde 2007 el PIB por habitante español ha crecido menos (8,7~\% frente al 12,6~\%): España crece más en total, pero no converge en renta por habitante. Como solo la productividad eleva la renta por habitante a largo plazo, este es el reto de fondo que recorre el resto del dictamen.

\textit{En suma, parece tratarse de una economía avanzada en una expansión madura, guiada por la demanda interna y sostenida por la población, con fortalezas (superávit exterior, desapalancamiento privado, banca saneada) y debilidades (productividad, paro estructural, deuda pública) que se desarrollan a continuación.}
\end{examen}

\begin{complemento}{1.1 · Por qué se contesta así: los 3 pasos y las fases}
\que{Qué se aplica.} Los 3 pasos del perfil cíclico (guía, I, gráfico «Los 3 pasos»): fases con el PIB y la brecha (I.2-I.3), motor por la demanda y la oferta (I.4-I.5) y diagnóstico anticipado (I.6). La frase inicial es la «frase ejecutiva» del FMI y del Banco de España (I.2): \emph{dónde está, qué la sostiene, qué la frena}; se escribe al final, en un hueco que se deja al empezar. Como referencia, el FMI resumía así su evaluación de 2026: la economía española «siguió creciendo más que la zona del euro pese al impacto adverso de la guerra en Oriente Medio» \parencite{FMI2026ArticuloIV}.

\que{Las fases, en un gráfico.} El cuadernillo no trae la brecha dibujada; el gráfico~\ref{fig:s01} la pone junto al PIB para ver las 4 fases de la guía (recuperación, expansión, desaceleración, recesión) en los datos.

\begin{figuracaja}
\caption{España: crecimiento del PIB real y brecha de producción, 2007-2027}\label{fig:s01}
\grafico[0.98\linewidth]{s01-fases-ciclo}
\fuente{Cuadro~1 del cuadernillo (Eurostat; brecha y previsiones de 2026-2027: Comisión Europea, AMECO, primavera de 2026). Qué hay que ver: la brecha toca fondo en 2013 aunque el PIB ya caía menos (\emph{pasar de caer un 2,9~\% a un 1,4~\% no es recuperación}); se cierra en 2017 y vuelve a ser positiva desde 2022; en 2026-2027 el PIB crecería menos que el potencial y la brecha bajaría: desaceleración.}
\end{figuracaja}

\begin{tablacaja}
\centering\scriptsize
\caption{Las fases en cifras (medias del periodo; brecha y paro, en el último año)}\label{tab:fases}
\begin{tabular}{>{\raggedright\arraybackslash}p{3.4cm}*{5}{>{\centering\arraybackslash}p{1.75cm}}}
\toprule
\rowcolor{tceeazulpalido}
\textbf{Fase} & \textbf{PIB (\%)} & \textbf{Demanda nacional (pp)} & \textbf{Demanda externa (pp)} & \textbf{Brecha (\% del potencial)} & \textbf{Paro (\%)} \\
\midrule
Final de la expansión, 2007-2008 & 2,2 & 1,7 & 0,5 & 1,1 & 11,3 \\
Doble recesión, 2009-2013 & $-1{,}7$ & $-3{,}5$ & 1,8 & $-9{,}6$ & 26,1 \\
Recuperación y expansión, 2014-2019 & 2,6 & 2,6 & 0,0 & 3,1 & 14,1 \\
Pandemia, 2020 & $-10{,}9$ & $-8{,}8$ & $-2{,}2$ & $-8{,}5$ & 15,5 \\
Rebote, 2021-2023 & 5,2 & 4,2 & 1,0 & 1,0 & 12,2 \\
Expansión, 2024-2025 & 3,2 & 3,4 & $-0{,}3$ & 2,0 & 10,5 \\
\bottomrule
\end{tabular}
\fuente{Cálculo propio con el cuadro~1 y el cuadro~3 (\texttt{calculos.txt}). Qué hay que ver: en las recesiones la demanda externa suma y la nacional resta; la recuperación de 2014-2019 fue de demanda nacional, con la externa neutra en promedio.}
\end{tablacaja}

\que{Errores típicos.} Contar la serie año a año («en 2010 crece un 0,1~\%, en 2011 cae un 0,6~\%…»): la guía pide la visión de conjunto, con lo atípico como excepción (I.2). Llamar «recuperación» a 2013, cuando el PIB todavía cae. Decir que la demanda interna «crece un 3,4~\%» en 2025: son 3,4 \emph{puntos} de los 2,6 que crece el PIB (I.4). Desarrollar aquí la balanza de pagos o la inflación, que tienen su pregunta: en el perfil cíclico la demanda externa se despacha en 1 o 2 frases (nota al opositor de I.4).
\end{complemento}

\begin{complemento}{1.2 · Los cálculos de la respuesta}
\que{Tasa anualizada} (guía, A.3.1): $(1+g)^4-1$; con $g=0{,}63~\%$, un 2,5~\%, y con $g=0{,}68~\%$, un 2,8~\%. Sin calculadora, $4g$ da 2,5 y 2,7: basta.

\que{Crecimiento medio de 2026 con los datos trimestrales.} Encadenando las tasas del gráfico~1 (base: cuarto trimestre de 2024 = 100), la media de 2025 es 101,5 y el segundo trimestre de 2026 está en 103,8. Si el PIB no creciera nada en la segunda mitad, 2026 crecería un 2,2~\% de media (solo por el «arrastre»); con un 0,4~\% trimestral, un 2,5~\%. La previsión de la Comisión (2,4~\%) supone, por tanto, algo menos del 0,4~\% por trimestre: una desaceleración suave, no un frenazo.

\que{Potencial que implica la brecha} (guía, I.3): $\gamma_{Y^p}\approx\gamma_Y-\Delta\mathrm{OG}$. En 2025, $2{,}6-(2{,}0-1{,}9)=2{,}5$, frente al 2,7 del cuadro; en 2024, $3{,}7-0{,}9=2{,}8$, frente al 2,5. Las diferencias se deben a que la brecha y el potencial son de AMECO (primavera de 2026), anterior a la última revisión del PIB del INE (2025: 2,8~\% en AMECO y 2,6~\% en el cuadro): \uline{si 2 fuentes no cuadran, se dice}, sin construir el análisis sobre décimas.

\que{Contribuciones} (guía, I.4): 2025, $3{,}4+(-0{,}8)=2{,}6$; 2009, $-6{,}5+2{,}8=-3{,}7\approx-3{,}8$ (los volúmenes encadenados no suman exacto).

\que{PIB por habitante} (gráfico~2): $118{,}7/107{,}9=1{,}100$ (+10,0~\%) y $108{,}7/103{,}6=1{,}050$ (+5,0~\%); población acumulada 2019-2025 (cuadro~3): $1{,}008\times1{,}005\times1{,}000\times1{,}009\times1{,}012\times1{,}011\times1{,}010\approx1{,}048$.
\end{complemento}

\begin{complemento}{1.3 · El PIB por habitante, descompuesto}
\que{La identidad} (guía, II.1, gráfico «PIB por persona en edad de trabajar»): en logaritmos,
\[
\gamma_Y=\underbrace{\gamma_N}_{\text{población}}+\underbrace{\gamma_{P_{15\text{-}64}/N}}_{\text{peso de la edad de trabajar}}+\underbrace{\gamma_{L/P_{15\text{-}64}}}_{\text{tasa de ocupación}}+\underbrace{\gamma_{H/L}}_{\text{horas por ocupado}}+\underbrace{\gamma_{Y/H}}_{\text{PIB por hora}} ,
\]
y el PIB por habitante es todo menos el primer término. El cuadro~2~bis la trae ya calculada; el gráfico~\ref{fig:s02} la dibuja por periodos.

\begin{figuracaja}
\caption{Descomposición del crecimiento del PIB por periodos (puntos al año)}\label{fig:s02}
\grafico[0.8\linewidth]{s02-descomposicion}
\fuente{Cuadro~2~bis del cuadernillo (Eurostat). Qué hay que ver: en la crisis la productividad por hora «sube» porque se destruye el empleo menos productivo (efecto composición, guía I.5); en 2013-2019 manda la ocupación; desde 2019 pesan la población y la ocupación, las horas por ocupado restan y la productividad por hora apenas suma.}
\end{figuracaja}

\que{Cómo leerla.} 3 periodos, 3 patrones: 2007-2013, ajuste por cantidades (la ocupación resta 2,9 puntos al año) con una productividad que sube de forma espuria; 2013-2019, recuperación del empleo con la población casi estancada; 2019-2025, \emph{más gente y menos horas}. La caída de las horas por ocupado ($-0{,}46$) puede deberse a más empleo a tiempo parcial, a la composición del empleo nuevo o a cambios de jornada: con el cuadernillo no se puede separar, así que se dice con condicional.

\que{Alternativa de redacción} (más corta, para el examen de 1 h 45 min): «Desde 2019 el PIB por habitante crece la mitad que el PIB (5,0 frente a 10,0~\%) porque la población aumenta casi un 5~\% por la inmigración; el crecimiento es extensivo: de 1,6 puntos anuales, 0,8 los pone la población y solo 0,3 la productividad por hora».

\que{Qué añadir con más tiempo.} La renta nacional frente al PIB (guía, I.3): las rentas primarias son casi nulas (+0,2~\% del PIB en 2025, cuadro~7), así que la RNB y el PIB apenas difieren, a diferencia de 2007-2008 ($-2{,}5$ y $-2{,}8~\%$), cuando España pagaba al exterior los intereses de su deuda. Y los indicadores adelantados (PMI, confianza), que el cuadernillo no trae: se dice que faltan.

\que{Qué recortar en el examen.} Quedarse con la frase de diagnóstico, las 5 fases en una línea cada una, 1 frase de demanda, 1 de oferta, el dato de 2026 con la previsión, 2 frases de brecha y la descomposición en 1,59 = 0,79 + 0,95 $-$ 0,46 + 0,33. Unas 500 palabras.
\end{complemento}
